\section{Evaluation protocol}

The empirical program now has two layers. Checkpoint 03 is a frozen paper-scale evaluation of the simpler pooled forecast-CV selector and remains a baseline. The central dynamic extension evaluates tracked local-minimum branches and alternative rules for converting branch history into the smoothness used at the current forecast origin.

\subsection{Layer 1: pooled forecast-CV baseline}

The frozen pooled baseline fixes
\[
d=2,\qquad L=60,\qquad h\in\{1,3,6,12\}.
\]
It compares horizon-matched pooled forecast-CV with one-step forecast-CV, ordinary CV, GCV, and AICc. BIC remains an exploratory diagnostic because its globally searched score was pinned to the left finite boundary in the refinement stage.

The pooled study uses controlled latent mechanisms, iid Gaussian, AR(1), and finite-variance Student-$t$ noise, fixed outer origins, and untouched future blocks. Its design was frozen before the paper-scale Checkpoint 03 run and is not retuned after observing those results.

\subsection{Layer 2: dynamic tracked-minimum evaluation}

The central dynamic experiment operates on a chronological sequence of Validation-1 forecast-loss surfaces. At each origin $t$, all relevant local minima
\[
\mathcal M_t=\{S_{1,t},\ldots,S_{K_t,t}\}
\]
are recovered. Local minima are linked across adjacent origins by one-to-one matching under
\[
|S_{j,t}-S_{j,t-1}|\le\varepsilon.
\]
The resulting branch history is stored as
\[
V_j
=
\{(S_{j,t},\ell^{(1)}_{j,t},\ell^{(2)}_{j,t})\}_{t\in\mathcal T_j},
\]
where $\ell^{(1)}$ is Validation-1 forecast loss and $\ell^{(2)}$ is Validation-2 forecast loss after refitting through Validation 1.

The dynamic protocol separates branch selection from final smoothness selection. A branch rule
\[
\widehat j_T=\psi(V_1,\ldots,V_J)
\]
uses historical branch support and Validation-2 performance only. Conditional on the selected branch, a rule
\[
\widehat S_T=\phi(V_{\widehat j_T})
\]
produces the smoothness used for the current forecast.

The first comparison set for $\phi$ contains:
\begin{enumerate}
\item the newest tracked local minimum;
\item a recent arithmetic mean;
\item a recent median;
\item a Validation-2-loss-weighted mean;
\item a recency-and-Validation-2 weighted mean.
\end{enumerate}
A direct forecast of the smoothness trajectory is reserved as a later extension unless its specification can be frozen without using outer-test information.

\subsection{Selection, refitting, and information boundaries}

Validation selects hyperparameters or branch summaries; it does not select a historical trend fit to carry into the test. For every candidate branch rule, after $\widehat S_T$ is obtained, all temporary Validation-1 and Validation-2 fits are discarded. The trend is freshly estimated using the most recent full window available at the outer origin:
\[
x_T=y_{T-L+1:T},
\qquad
\widehat\tau_T
=
H_{\lambda(\widehat S_T)}x_T.
\]
The operational forecast is
\[
\widehat y_{T+1:T+h\mid T}
=
G_{d,h}\widehat\tau_T.
\]
The outer future block is revealed only for scoring. At a later outer origin, observations that have become available enter the historical information set and the complete branch-selection, smoothness-selection, refit, and forecast cycle is repeated.

Thus the procedure may average historical forecast losses or historical smoothness values when a declared rule requires it, but it never averages historical trend fits.

\subsection{Fair comparison of branch-to-smoothness rules}

All $\phi$ rules must be evaluated on exactly the same recovered branch histories, the same branch selector $\psi$, and the same untouched outer test blocks. Hyperparameters such as the continuation radius $\varepsilon$, recent-history length $K$, recency factor $\rho$, and weighting stabilizer $\delta$ must be fixed from development data before final test scoring.

The dynamic experiment should report forecast MSFE and MAE, selected smoothness, branch support, branch switching or disappearance rates, and the distance between each feasible rule and a simulation-only latent forecast oracle where available. The pooled forecast-CV selector remains a direct baseline.

\subsection{Role of the numerical-methods companion paper}

The forecasting paper assumes that the relevant local minima on each surface and their cross-origin correspondences can be supplied reliably. Numerical discovery of those minima, endpoint handling, adaptive subdivision, derivative-root refinement, and the branch-correspondence problem belong to the companion numerical-methods paper. The forecasting paper evaluates what information in the resulting branch matrices is useful for selecting the current smoothness.

\subsection{Public real-data evaluation}

The final paper should complement simulation with a heterogeneous public panel under the same chronology. No series-specific branch rule may be chosen from its test performance. If macroeconomic histories are used, historical backtests require real-time vintages rather than revised current histories.

\subsection{Freeze rules}

The frozen Checkpoint 03 pooled design remains unchanged. The dynamic branch experiment receives its own freeze point: the branch matcher, $\psi$, candidate $\phi$ rules, $\varepsilon$, $K$, $\rho$, $\delta$, horizons, origins, and primary metrics must be recorded before the final dynamic paper-scale run. Unfavorable final results do not justify changing those choices.
